\documentclass[11pt]{article}
\usepackage{hyperref}
\usepackage{enumitem}
\usepackage{cmap}
\usepackage[utf8]{inputenc}
\usepackage[T1]{fontenc}
\usepackage[left=1.15cm,top=1.35cm,right=1.15cm,bottom=0.9cm]{geometry}
\IfFileExists{glyphtounicode.tex}{\input{glyphtounicode}\pdfgentounicode=1}{}
\hypersetup{
    hidelinks,
    pdfauthor={Kevin Bravo},
    pdftitle={Kevin Bravo - Frontend Developer - React, Astro, SvelteKit y TypeScript},
    pdfsubject={CV para el proyecto de landing de Spider-Man Brand New Day con InfoJobs y Midudev},
    pdfkeywords={Kevin Bravo, Frontend Developer, Programador Frontend, React, Astro, SvelteKit, Next.js, TypeScript, JavaScript, HTML, CSS, Tailwind CSS, responsive design, rendimiento web, SSR, SEO, Cloudflare, landing page, InfoJobs, Midudev}
}
\setlength{\parindent}{0pt}
\setlist[itemize]{leftmargin=*, noitemsep, topsep=2pt, partopsep=0pt, parsep=2pt, label={\textbf{-}}}

\begin{document}

\begin{center}
    {\LARGE \textbf{Kevin Bravo}}\\
    Frontend Developer | React, Astro, SvelteKit, TypeScript \\
    \hrulefill
\end{center}

\begin{center}
    \href{mailto:0bravokevin@gmail.com}{0bravokevin@gmail.com}
    \textbullet\
    \href{https://kevinbravo.com}{kevinbravo.com}
    \textbullet\
    \href{https://github.com/0bkevin}{github.com/0bkevin}
    \textbullet\
    \href{https://www.linkedin.com/in/0bkevin/}{linkedin.com/in/0bkevin}
\end{center}

\vspace{4pt}
\begin{center}
    \textbf{Perfil profesional}
\end{center}
\vspace{-4pt}
Desarrollador frontend y full-stack con 5 años de experiencia creando productos web, interfaces responsivas, dashboards y plataformas públicas con TypeScript, React, SvelteKit y Astro. He liderado la migración de una aplicación activa de React a SvelteKit, mejorando su arquitectura frontend, SSR, mantenibilidad y rendimiento. Desarrollo productos completos desde el diseño de la experiencia y la implementación visual hasta las APIs, el despliegue y la optimización en producción. Trabajo bien con plazos cortos, requisitos cambiantes y equipos multidisciplinarios, con atención al detalle y foco en una navegación rápida y clara.

\vspace{8pt}
\begin{center}
    \textbf{Habilidades técnicas}
\end{center}
\vspace{-4pt}
\textbf{Frontend:} React, Svelte 5, SvelteKit, Astro, Next.js, TypeScript, JavaScript, HTML5, CSS3, Tailwind CSS, interfaces responsivas, componentización \\
\textbf{Rendimiento y calidad web:} SSR, SEO técnico, metadatos, sitemap, caché HTTP, accesibilidad, optimización de carga, diseño mobile-first, pruebas funcionales \\
\textbf{Backend e integraciones:} Node.js, REST APIs, Python, Django, PostgreSQL, Drizzle ORM, Prisma, integraciones con terceros \\
\textbf{Entrega:} Git, GitHub, Cloudflare Workers, Azure, Docker, debugging, despliegues, revisión de releases, trabajo remoto

\vspace{8pt}
\begin{center}
    \textbf{Proyecto destacado}
\end{center}
\vspace{-4pt}
\textbf{RegistroCiberVE} \hfill \textit{2026 -- Presente} \\
\textit{Svelte 5, SvelteKit 2, TypeScript, Tailwind CSS, ECharts, PostgreSQL, Cloudflare Workers}
\begin{itemize}
    \item \textbf{Diseño y desarrollo un registro público bilingüe de incidentes de ciberseguridad}, con landing, directorios de incidentes, atacantes y víctimas, búsqueda, filtros, perfiles detallados y visualizaciones responsivas.
    \item \textbf{Implementé una experiencia preparada para distribución pública}, con SSR, SEO, rutas canónicas, sitemap dinámico, imágenes para compartir, QR y contenido en español e inglés.
    \item \textbf{Desplegué la aplicación en Cloudflare Workers} con PostgreSQL en Azure mediante Hyperdrive, caché, Turnstile, WAF y controles de abuso para mantener tiempos de respuesta y disponibilidad.
\end{itemize}

\vspace{8pt}
\begin{center}
    \textbf{Experiencia}
\end{center}
\vspace{-4pt}

\textbf{Proyecto Educa} \\
\textbf{Chief Technology Officer} \hfill \textit{Dic 2025 -- Presente}
\begin{itemize}
    \item \textbf{Dirijo el desarrollo de una plataforma de e-learning} para impartir y gestionar con escuelas cursos dirigidos a estudiantes de 12 a 18 años.
    \item \textbf{Soy responsable de la experiencia de producto y la entrega técnica}, desde las interfaces y la arquitectura hasta los releases y el soporte en producción.
\end{itemize}

\vspace{8pt}
\textbf{Free2Z} \\
\textbf{Principal Software Engineer} \hfill \textit{Abr 2025 -- Abr 2026}
\begin{itemize}
    \item \textbf{Lideré la migración de un producto en producción de React a SvelteKit}, mejorando la arquitectura frontend, SSR, mantenibilidad, manejo de estado y rendimiento sin interrumpir el servicio.
    \item \textbf{Entregué funcionalidades orientadas a usuarios} integradas con APIs en Python y Django, video en vivo con Dyte y una suite CMS con editor Markdown y motor de blog.
\end{itemize}

\vspace{8pt}
\textbf{U.S. Department of State (via AVAA)} \\
\textbf{Product Engineer | Billingua Talent} \hfill \textit{Nov 2024 -- Ago 2025}
\begin{itemize}
    \item \textbf{Construí Billingua Talent como único ingeniero}, desde el modelo de datos hasta perfiles, dashboards, flujos de registro, interfaces por rol, despliegue y cada release.
    \item \textbf{Probé el producto con 50 usuarios y 5 empresas} y convertí el feedback del piloto en mejoras de navegación, flujos y reportes.
\end{itemize}

\textbf{Freelance} \\
\textbf{Software Engineer} \hfill \textit{Sep 2024 -- Dic 2024}
\begin{itemize}
    \item \textbf{Desarrollé sitios y productos web para clientes} con Astro, TypeScript y Tailwind CSS, cuidando diseño responsivo, formularios, estructura de contenido, rendimiento y despliegue.
    \item \textbf{Me encargué del proceso completo}, desde entender el objetivo y proponer la interfaz hasta implementar, publicar, revisar y ajustar el resultado con el cliente.
\end{itemize}

\vspace{8pt}
\textbf{Asociación Venezolano Americana de Amistad (AVAA)} \\
\textbf{Software Developer} \hfill \textit{Oct 2022 -- Jul 2024}
\begin{itemize}
    \item \textbf{Propuse, diseñé y construí SEP}, una plataforma web que hoy atiende a cerca de 300 becarios y equipos internos en tres capítulos nacionales.
    \item \textbf{Desarrollé SEA para EducationUSA}, con historiales, servicios, pagos, recibos, dashboards y reportes en una interfaz centralizada para los asesores.
\end{itemize}

\vspace{8pt}
\begin{center}
    \textbf{Colaboración y comunidad}
\end{center}
\vspace{-4pt}
\textbf{Dev3pack, Organizador del hub del Global Hackathon} \hfill \textit{Abr 2026 -- May 2026} \\
Organicé el hub de Caracas de un hackathon global de Web3 e IA, coordinando sede, agenda, aliados, mentores y participantes; también asesoré a los equipos para definir alcance, resolver decisiones técnicas y entregar a tiempo. \\
\textbf{Technovation Girls Venezuela, Mentor de programación} \hfill \textit{2026} \\
Guié a un equipo estudiantil en la definición del problema, el prototipado y el desarrollo de una aplicación móvil con MIT App Inventor.

\vspace{8pt}
\begin{center}
    \textbf{Educación e idiomas}
\end{center}
\vspace{-4pt}
\textbf{Instituto de Estudios Superiores de Administración (IESA)} \hfill \textit{En curso} \\
Maestría en Administración Pública (MPA) \\
\textbf{Idiomas:} Español (nativo), Inglés (C1)

\end{document}
